\section{综合案例与练习卷}

一本操作系统书要真正有用，必须让读者能把章节串起来。本节给出几个贯穿案例。它们不是实践卷工程，但可以作为理论卷的验收题：读者应能写出状态机、关键对象、不变量、失败路径和测试证据。

\subsection{案例一：shell 执行 \code{cmd > out.txt}}

这个案例串起进程、fd、路径、exec、调度和文件写入。

\begin{lstlisting}
shell:
  fd = open("out.txt", O_CREAT | O_TRUNC | O_WRONLY)
  pid = fork()
  if pid == 0:
    dup2(fd, STDOUT_FILENO)
    close(fd)
    execve("cmd", argv, envp)
  close(fd)
  waitpid(pid)
\end{lstlisting}

要求回答：

\begin{enumerate}
\tightlist
\item fork 后父子进程的地址空间、fd 表和 open file description 如何关联？
\item dup2 后 stdout 指向哪个 file object？
\item exec 成功后哪些 fd 保留，哪些会因为 \code{CLOEXEC} 关闭？
\item cmd 写 stdout 时，数据经过哪些内核对象？
\item waitpid 为什么需要 zombie 状态？
\end{enumerate}

\subsection{案例二：HTTP 服务端处理一个请求}

这个案例串起 socket、epoll、线程池、内存分配、文件读取和 send。

\begin{lstlisting}
event loop:
  events = epoll_wait()
  for event in events:
    if event.fd == listen_fd:
      accept connections
    else:
      read request bytes
      parse frame
      dispatch worker

worker:
  open file
  read through page cache
  send response with short-write loop
\end{lstlisting}

要求回答：

\begin{enumerate}
\tightlist
\item listen socket 和 accepted socket 是不是同一个对象？
\item 为什么一次 recv 不等于一个完整 HTTP 请求？
\item 文件热缓存和冷缓存下，read 路径差异在哪里？
\item send 返回 \code{EAGAIN} 时，用户态状态机如何保存剩余数据？
\item 服务端关闭连接时，哪些等待队列要被唤醒？
\end{enumerate}

\subsection{案例三：数据库提交一条记录}

这个案例串起 page cache、日志、fsync、rename、块层 flush 和崩溃恢复。

\begin{lstlisting}
transaction:
  append WAL intent
  fsync WAL
  update data page
  append commit record
  fsync WAL
  later checkpoint data pages
\end{lstlisting}

要求回答：

\begin{enumerate}
\tightlist
\item 为什么 intent 记录要先于数据页持久化？
\item commit record 写入前崩溃，恢复应如何处理？
\item commit record 写入后数据页未写完，恢复为什么可以 redo？
\item checksum 如何帮助识别 torn write？
\item checkpoint 为什么不能截断尚未安全覆盖的数据所需日志？
\end{enumerate}

\subsection{案例四：系统变慢的诊断}

现象：请求 p99 从 80ms 升到 3s，CPU 使用率不高。

要求写出诊断路径：

\begin{lstlisting}
1. collect thread states
2. split on-CPU and off-CPU time
3. inspect syscall latency
4. inspect page fault and reclaim
5. inspect block I/O queue and flush
6. inspect cgroup throttling
7. form evidence chain
\end{lstlisting}

合格报告必须说明排除了什么。比如 CPU profiler 没有热点，只能说明 CPU 不是主因，不能自动证明磁盘是主因；还需要 syscall trace、block trace 或 dirty page 证据。

\subsection{期末式综合题}

\begin{enumerate}
\item 解释从定时器中断到任务抢占的完整路径。要求包含 trap frame、\code{need\_reschedule}、\code{preempt\_count}、runqueue 和上下文切换。
\item 解释从用户态 \code{write(fd, buf, len)} 到 page cache dirty page 的完整路径。要求包含 fd、file、inode、用户指针复制、短写和错误码。
\item 解释从 \code{mmap} 文件页第一次读取到 page fault 完成的完整路径。要求包含 VMA、PTE、page cache、块 I/O 和 TLB。
\item 解释一个设备 DMA 写入内存后，CPU 如何安全读取数据。要求包含 DMA map、cache invalidate、completion 和 buffer 生命周期。
\item 解释容器隔离为什么不是虚拟机。要求包含 namespace、cgroup、capability、seccomp、共享内核和逃逸风险。
\item 设计一个崩溃一致的多文件发布协议。要求包含临时文件、fsync、rename、父目录 fsync、manifest、checksum 和恢复规则。
\end{enumerate}

\begin{rubricbox}
综合题评分重点不是背 API，而是能不能把对象、状态、权限、等待、提交和证据串成闭环。只写正常路径不给满分；必须写错误路径和边界。
\end{rubricbox}

\subsection{参考答题框架}

综合题不要求背 Linux 函数名，而要求结构完整。推荐按下面框架答题：

\begin{lstlisting}
Objects:
  涉及哪些用户态、内核、硬件对象

Normal path:
  正常路径逐步推进，每步改变什么状态

Permission:
  哪些地方做权限或边界检查

Waiting:
  哪些地方可能阻塞，等待队列在哪里

Failure:
  每个失败点如何返回、回滚或提交部分结果

Observation:
  用什么日志、trace、计数器或测试证明结论
\end{lstlisting}

这个框架能防止答案只写“调用 A 然后调用 B”。操作系统题目的本质是对象和状态，不是函数名接龙。

\subsection{样例解析：从 \code{write} 到持久化}

一个完整答案应类似下面这样展开：

\begin{enumerate}
\tightlist
\item 用户态调用 \code{write(fd, buf, len)}，fd 是进程局部整数。
\item 内核在 fd table 找到 file object，检查是否允许写。
\item 内核用 \code{copy\_from\_user} 分段复制用户 buffer，失败返回 \code{EFAULT}。
\item 普通文件写入 page cache，相关页变 dirty，inode size 可能更新。
\item \code{write} 返回写入字节数；这不是持久化边界。
\item 后台 writeback 或 \code{fsync} 构造 bio，经块层和驱动提交到设备。
\item 设备 completion 后，页可从 writeback 变 clean。
\item 若要业务提交，还需要目录、rename、manifest 或日志协议。
\end{enumerate}

错误路径也必须写：fd 无效返回 \code{EBADF}；用户指针坏返回 \code{EFAULT}；空间不足可能返回 \code{ENOSPC}；设备错误可能在 \code{fsync} 或 close 时报出。只写正常路径不合格。

\subsection{样例解析：线程卡死}

先不要猜锁。按顺序排查：

\begin{enumerate}
\tightlist
\item 线程是 RUNNABLE 还是 SLEEPING？
\item 若 RUNNABLE，看 runqueue、优先级、CPU 配额和抢占。
\item 若 SLEEPING，看 wait channel：锁、I/O、page fault、futex、timer 还是 waitpid。
\item 若等锁，找持锁者和锁顺序图。
\item 若等 I/O，看块层队列、设备完成和 dirty writeback。
\item 若等 futex/条件变量，看等待谓词和关闭路径是否唤醒。
\item 最后用 trace 或日志证明等待原因，而不是只看现象。
\end{enumerate}

\subsection{全书复习清单}

\begin{itemize}
\tightlist
\item 能从硬件解释特权级、页表、中断、DMA 和启动顺序。
\item 能从状态机解释进程、线程、等待、唤醒和退出。
\item 能从边界解释系统调用、用户指针、错误码和权限。
\item 能从缓存解释 page cache、mmap、writeback 和 fsync。
\item 能从队列解释块层、设备驱动、socket 和 epoll。
\item 能从恢复规则解释日志、manifest、checksum 和幂等。
\item 能从证据链解释性能、安全、panic 和资源泄漏。
\end{itemize}

\subsection{案例五：从按键中断到终端程序读到字符}

这个案例把硬件、中断、驱动、TTY、等待队列和用户态读串起来。一个键盘或串口设备产生输入事件，设备控制器触发中断；驱动读取设备寄存器，把扫描码或字节放入内核缓冲；TTY 层处理行规程、回显和特殊字符；正在 \code{read} 的 shell 被唤醒，复制数据到用户 buffer。

\begin{lstlisting}
keyboard interrupt:
  save trap frame
  driver reads scancode
  translate scancode to key event
  tty line discipline updates buffer
  wake read_waitq

shell read:
  if tty buffer empty:
    sleep(read_waitq)
  copy chars to user buffer
  return byte count
\end{lstlisting}

要求回答：

\begin{enumerate}
\tightlist
\item 中断处理函数为什么不能等待用户程序消费字符？
\item TTY 缓冲区满时，驱动和上层分别有哪些选择？
\item \code{Ctrl-C} 为什么不是普通字符那么简单？
\item shell 阻塞在 \code{read} 时，调度器看到的线程状态是什么？
\item 如果用户传入坏 buffer，错误应该在哪里被发现？
\end{enumerate}

\subsection{案例六：fork 后多线程程序的陷阱}

多线程程序调用 \code{fork} 时，子进程通常只复制调用 fork 的那个线程。其他线程不在子进程继续运行，但它们持有的锁状态可能被复制进来。如果子进程在 \code{exec} 前调用复杂库函数，就可能等待一个永远没人释放的锁。

\begin{lstlisting}
thread A:
  lock(malloc_lock)
  ... interrupted here ...

thread B:
  pid = fork()

child:
  // only thread B exists, but malloc_lock may look locked
  printf("hello")  // may call malloc and deadlock
  execve(...)
\end{lstlisting}

要求回答：

\begin{enumerate}
\tightlist
\item fork 复制了哪些进程状态，哪些线程不会继续存在？
\item 为什么 fork 后 exec 前只能调用 async-signal-safe 的少量函数？
\item \code{pthread\_atfork} 试图解决什么问题？
\item COW 如何降低 fork 成本，但不能解决锁状态复制问题？
\item shell 的 fork/exec 模型为什么在单线程 shell 中更简单？
\end{enumerate}

\subsection{案例七：内存压力导致网络服务超时}

现象：服务端 CPU 不高，网络连接没有明显丢包，但请求大量超时。诊断发现 p99 卡在 \code{send} 或 \code{write} 附近。可能路径是：请求处理需要分配 buffer；内存压力触发回收；回收要写回脏页；块设备 flush 慢；线程 sleeping；socket 发送迟迟得不到调度。

\begin{lstlisting}
request handler:
  allocate response buffer
  -> reclaim
  -> dirty page writeback
  -> block device wait
  -> scheduler sleeps current
  -> client timeout
\end{lstlisting}

要求回答：

\begin{enumerate}
\tightlist
\item 为什么 CPU 低不代表服务健康？
\item 哪些指标能区分锁等待、I/O 等待和 cgroup throttling？
\item dirty page 背压和 socket 背压有什么相似处？
\item 如果 OOM killer 触发，应该在日志中看到哪些字段？
\item 如何设计限流，让服务在内存压力下早失败而不是拖到超时？
\end{enumerate}

\subsection{案例八：容器内服务被限制后无法写日志}

一个容器服务启动正常，但运行时写日志返回只读文件系统错误。排查发现根文件系统以只读方式挂载，只有 \code{/var/log/app} 应该挂可写卷，但实际没有挂载。这个问题串起 mount namespace、VFS、路径解析、权限和观测。

\begin{lstlisting}
open("/var/log/app/current.log", O_CREAT|O_APPEND|O_WRONLY)
  -> path lookup in container mount namespace
  -> parent directory resolves to readonly mount
  -> VFS denies write
  -> return -EROFS
\end{lstlisting}

要求回答：

\begin{enumerate}
\tightlist
\item 为什么宿主上路径存在，不代表容器 namespace 中可见？
\item \code{EACCES}、\code{EROFS}、\code{ENOENT} 分别说明什么？
\item 这个失败应记录哪些安全/运维字段？
\item capability 能不能绕过只读挂载？为什么不应依赖它？
\item 测试镜像时如何证明只读根文件系统和可写日志卷都按预期工作？
\end{enumerate}

\subsection{算法推导题：三种页面置换策略}

给定页面访问序列：

\begin{lstlisting}
7, 0, 1, 2, 0, 3, 0, 4, 2, 3, 0, 3, 2
\end{lstlisting}

物理页框数量为 3。分别手算 FIFO、LRU 和 Clock 的缺页次数，并回答：

\begin{enumerate}
\tightlist
\item 为什么 FIFO 可能出现 Belady anomaly？
\item 精确 LRU 的维护成本在哪里？
\item Clock 的 referenced bit 何时清零，何时给第二次机会？
\item 脏页和干净页在回收成本上有什么差异？
\item 文件页和匿名页在回收路径上有什么差异？
\end{enumerate}

\subsection{算法推导题：调度延迟和响应时间}

给定三个任务：

\begin{longtable}{p{0.16\linewidth}p{0.22\linewidth}p{0.22\linewidth}p{0.26\linewidth}}
\toprule
任务 & 到达时间 & CPU 总需求 & 行为 \\
\midrule
A & 0 & 24 tick & CPU 密集 \\
B & 2 & 6 tick & 每运行 1 tick 后 I/O 等待 4 tick \\
C & 4 & 3 tick & 短任务 \\
\bottomrule
\end{longtable}

分别分析 FCFS、RR（时间片 2）、MLFQ（两级队列）下的完成时间、平均等待时间和 B 的最大响应延迟。要求说明：

\begin{enumerate}
\tightlist
\item 为什么 FCFS 会产生 convoy effect？
\item 时间片变小如何影响响应和上下文切换成本？
\item MLFQ 为什么倾向于保留 B 的高优先级？
\item 若 A 持有 B 需要的锁，单纯提高 B 优先级是否足够？
\item 如何用 trace 证明响应延迟来自调度而不是 I/O？
\end{enumerate}

\subsection{算法推导题：日志恢复}

给定如下日志片段，其中每条记录都有 LSN、事务号、类型和 checksum。假设最后一条记录 checksum 损坏：

\begin{lstlisting}
LSN 10: TX1 INTENT write page P old=a new=b checksum ok
LSN 11: TX1 COMMIT checksum ok
LSN 12: TX2 INTENT write page Q old=x new=y checksum ok
LSN 13: TX2 COMMIT checksum bad
\end{lstlisting}

回答：

\begin{enumerate}
\tightlist
\item redo 日志恢复时 TX1 和 TX2 分别如何处理？
\item 如果是 undo 日志，记录内容和恢复规则会有什么不同？
\item 为什么 checksum 损坏后不能继续信任后续日志尾部？
\item checkpoint 能缩短哪些扫描，不能省略哪些安全检查？
\item manifest 发布协议如何用类似思想只采用完整版本？
\end{enumerate}

\subsection{设计题：最小教学内核的验收标准}

虽然本册暂不实现实践卷，但读者应能写出一个最小教学内核的验收标准。该内核不必追求性能，但必须证明机制边界清楚。

\begin{enumerate}
\item 启动：能从 bootloader 入口建立栈、清 BSS、输出早期日志、解析内存地图。
\item 中断：能注册时钟中断和一个简单设备中断，保存 trap frame 并返回。
\item 调度：至少支持两个内核线程抢占运行，阻塞线程不会被选中。
\item 同步：mutex 或 sleep queue 不丢唤醒，关闭路径唤醒所有等待者。
\item 虚存：能建立用户页表，用户态不能访问内核页，缺页错误可诊断。
\item 系统调用：能验证用户指针，返回区分 \code{EFAULT/EBADF/EINVAL} 的错误。
\item 文件抽象：至少有 fd 表和一个内存文件对象，支持 read/write/close 生命周期。
\item 观测：panic 输出 CPU、RIP、RSP、当前线程、最近系统调用和锁状态。
\end{enumerate}

\begin{rubricbox}
验收不是看功能数量，而是看每个功能是否有不变量、错误路径和测试证据。一个小而严谨的内核，比一个功能多但边界不清的内核更适合学习。
\end{rubricbox}

\subsection{答题评分细则}

为了把“会背概念”和“理解系统”区分开，综合题按五类给分：

\begin{longtable}{p{0.22\linewidth}p{0.58\linewidth}}
\toprule
维度 & 满分特征 \\
\midrule
对象 & 明确列出用户对象、内核对象、硬件对象和引用关系 \\
状态 & 写出正常状态转换、等待状态、关闭状态和错误状态 \\
边界 & 说明权限、用户指针、生命周期和持久化边界 \\
算法 & 给出关键算法或伪代码，并说明复杂度或取舍 \\
证据 & 给出测试、trace、日志、计数器或崩溃恢复验证办法 \\
\bottomrule
\end{longtable}

扣分规则也很明确：只写正常路径扣状态分；只写 API 名字扣对象分；不写错误码扣边界分；不写测试证据扣证据分；把 \code{write} 返回当作持久化提交、把 TCP 当作消息协议、把容器当作虚拟机，属于概念性错误。

\subsection{案例九：内核模块加载与 device mapper 创建}

这个案例串起权限、ELF/模块格式、符号解析、内核内存分配、设备模型和块层。加载 device mapper target 时，用户态通常通过 \code{modprobe} 或内核自动加载模块，随后通过 dm ioctl 创建映射设备。

\begin{lstlisting}
modprobe dm-snapshot
  -> finit_module()
  -> check CAP_SYS_MODULE and LSM
  -> copy module image from fd
  -> verify ELF sections and relocations
  -> resolve exported symbols
  -> allocate module memory
  -> apply relocations
  -> run module_init

dmsetup create vg-lv
  -> ioctl(/dev/mapper/control)
  -> parse target table
  -> create mapped device
  -> register gendisk/request path
\end{lstlisting}

要求回答：

\begin{enumerate}
\tightlist
\item 为什么加载模块需要比普通 open 更强的 capability？
\item 模块初始化失败时，已分配的 text/data、sysfs 节点和设备注册如何回滚？
\item device mapper 如何把上层 bio 映射到底层设备？
\item flush/FUA 请求穿过 dm target 时为什么不能丢失顺序语义？
\item 观测时应看哪些证据：dmesg、\filepath{/sys/module}、\filepath{/dev/mapper}、\filepath{/proc/devices} 还是 block trace？
\end{enumerate}

\subsection{案例十：容器从 clone 到 exec 的完整路径}

容器启动不是一个单独系统调用。它组合 clone/unshare、namespace、cgroup、mount、capability、seccomp、pivot root/chroot 和 exec。

\begin{lstlisting}
runtime:
  clone with CLONE_NEWUSER | CLONE_NEWPID | CLONE_NEWNS | CLONE_NEWNET
  write uid_map/gid_map
  configure cgroup v2 limits
  setup mount namespace and readonly rootfs
  drop bounding/ambient capabilities
  set no_new_privs
  install seccomp filter
  execve(container_init)
\end{lstlisting}

要求回答：

\begin{enumerate}
\tightlist
\item user namespace 为什么要先建立 uid/gid 映射？
\item pid namespace 中的 1 号进程退出时，容器内其他进程会怎样？
\item mount namespace 的只读根和可写 volume 如何同时存在？
\item seccomp 应在 exec 前还是后安装？各有什么取舍？
\item 如何证明容器不能 ptrace 宿主进程、不能 mount 宿主路径、不能突破内存限制？
\end{enumerate}

\subsection{案例十一：TCP 服务端从 socket 到 epoll}

这个案例把 fd、socket、网络栈、等待队列和 epoll ready list 串起来。

\begin{lstlisting}
server:
  listen_fd = socket()
  bind(listen_fd, addr)
  listen(listen_fd, backlog)
  epoll_ctl(ep, ADD, listen_fd, EPOLLIN)

event:
  packet arrives
  -> NIC/NAPI
  -> IP/TCP
  -> established socket receive queue
  -> sk_data_ready
  -> ep_poll_callback
  -> epoll_wait returns fd
\end{lstlisting}

要求回答：

\begin{enumerate}
\tightlist
\item listen socket 和 accepted socket 在内核对象上有什么区别？
\item SYN queue、accept queue、receive queue 分别堵塞时表现如何？
\item EPOLLET 下为什么必须读到 \code{EAGAIN}？
\item 慢客户端导致 send buffer 堵塞时，服务端状态机保存什么？
\item 如何用 ss、\filepath{/proc/net/tcp}、perf/ftrace 或 eBPF 证明瓶颈在网络栈还是应用队列？
\end{enumerate}

\subsection{案例十二：文件系统从 mount 到 fsck}

挂载文件系统时，内核读取 superblock，建立 VFS super\_block、root dentry 和 inode 缓存。崩溃后，journal replay 或 fsck 负责恢复结构一致性。

\begin{lstlisting}
mount /dev/sdb1 /mnt
  -> open block device
  -> read superblock
  -> validate block size and feature flags
  -> load journal if present
  -> create VFS super_block
  -> instantiate root dentry

after crash:
  -> journal replay committed transactions
  -> fsck scans inode/block bitmap if needed
\end{lstlisting}

要求回答：

\begin{enumerate}
\tightlist
\item mount 时为什么要检查 feature flags？
\item journal replay 和 fsck 分别解决什么问题？
\item ext4 ordered 模式为什么要求数据先于相关元数据 commit？
\item orphan inode list 的作用是什么？
\item fsck 如何发现 bitmap 标记空闲但 inode 引用该块的矛盾？
\end{enumerate}

\subsection{案例十三：从按键到字符设备到 read}

这个案例扩展前面的终端输入路径，要求读者把字符设备对象写清楚。

\begin{lstlisting}
keyboard IRQ
  -> input driver reads scancode
  -> input subsystem emits key event
  -> tty line discipline buffers char
  -> wake read waitqueue
  -> read(/dev/tty, buf, len)
  -> copy_to_user
\end{lstlisting}

要求回答：

\begin{enumerate}
\tightlist
\item 字符设备如何通过 major/minor 找到 \code{file\_operations}？
\item IRQ handler、input core、TTY line discipline 分别保存什么状态？
\item canonical mode 下为什么要等换行才返回？
\item \code{Ctrl-C} 如何变成 \code{SIGINT}，而不是普通字符？
\item 用户 buffer 坏时，已经从 TTY buffer 取出的字符如何处理？
\end{enumerate}

\subsection{案例十四：内存压力下的页面回收与 OOM 决策}

现象：服务进程不断分配内存，系统延迟升高，最后某个进程被杀。

\begin{lstlisting}
allocation slow path
  -> direct reclaim
  -> shrink_lruvec
  -> writeback dirty file pages
  -> shrink_slab
  -> compaction
  -> still below watermark
  -> OOM select victim
  -> kill and reap memory
\end{lstlisting}

要求回答：

\begin{enumerate}
\tightlist
\item active/inactive LRU 如何帮助选择回收页？
\item 干净文件页、脏文件页、匿名页、mlock 页分别如何处理？
\item cgroup OOM 和全局 OOM 的影响范围有何不同？
\item OOM killer 为什么不能随便杀内核线程或关键服务？
\item 如何用 \filepath{/proc/meminfo}、PSI、tracepoint 和 OOM 日志形成证据链？
\end{enumerate}

\subsection{案例十五：系统调用被信号中断后的重启路径}

阻塞系统调用遇到信号时，内核可能返回 \code{EINTR}，也可能在 handler 返回后重启。这个案例考察系统调用边界、等待队列、信号投递和用户可见语义。

\begin{lstlisting}
read(fd, buf, len)
  -> no data, sleep interruptible
  -> signal arrives
  -> wake task
  -> prepare signal frame
  -> handler runs in user mode
  -> sigreturn
  -> restart syscall or return EINTR
\end{lstlisting}

要求回答：

\begin{enumerate}
\tightlist
\item 哪些等待可以被信号中断，哪些不应被普通信号中断？
\item 已经完成部分读写时，为什么不能简单重启？
\item \code{SA\_RESTART} 改变了什么，不改变什么？
\item 内核内部的 \code{ERESTARTSYS} 一类返回值为什么不直接暴露给用户？
\item strace 中如何观察 restarted syscall？
\end{enumerate}
